\usepackage[
  margin=1.5cm, 
  top=1.8cm, 
  bottom=1.4cm,
  footskip=0.8cm,
  headsep=0.4cm,
]{geometry}
\usepackage{xcolor}
\usepackage{guitarchordschemes}
\usepackage{titlesec}
\usepackage[colorlinks=true, linkcolor=black]{hyperref}
\usepackage{subfiles}
\usepackage{fancyhdr}
\usepackage{tikz}
\usetikzlibrary{decorations.pathreplacing}

\usepackage{extramarks} % Provides \lastrightmark
\usepackage{ifthen}     % Provides conditional logic

\pagestyle{fancy}
\fancyhf{} % Clear all headers and footers

% Store titles without numbers
\renewcommand{\sectionmark}[1]{\markboth{#1}{}}
\renewcommand{\subsectionmark}[1]{\markright{#1}}

\fancyhead[R]{%
  \nouppercase{\leftmark}%
  \ifthenelse{\equal{\lastrightmark}{}}{}{ / \nouppercase{\lastrightmark}}%
}
\fancyfoot[R]{\thepage}

% 1. REMOVE THE HEADER LINE
\renewcommand{\headrulewidth}{0pt}

% ---------------------------------------------------------
% COLOR DEFINITIONS & TEXT LABEL SHORTCUTS
% ---------------------------------------------------------
\definecolor{rootcol}{RGB}{255,0,0}
\definecolor{thirdcol}{RGB}{0, 0, 255}
\definecolor{fifthcol}{RGB}{0, 175, 0}
\definecolor{seventhcol}{RGB}{255,0,255}
\definecolor{tensioncol}{RGB}{255,130,0} 

\newcommand{\R}{\textcolor{rootcol}{R}}
\newcommand{\third}[1]{\textcolor{thirdcol}{#1}}
\newcommand{\fifth}[1]{\textcolor{fifthcol}{#1}}
\newcommand{\seventh}[1]{\textcolor{seventhcol}{#1}}
\newcommand{\sixth}[1]{\textcolor{seventhcol}{#1}}
\newcommand{\tension}[1]{\textcolor{tensioncol}{#1}}

\newcommand{\s}{\ensuremath{\sharp}}
\renewcommand{\b}{\ensuremath{\flat}}
\newcommand{\bb}{\ensuremath{\flat\mkern-1.5mu\flat}}

% ---------------------------------------------------------
% DUAL NOTE LABEL MACROS (FIXED FOR TEXT MODE & TIKZ)
% ---------------------------------------------------------
% 1. Side-by-side slash label (using \hspace instead of math-only \mkern)
\newcommand{\duallabel}[2]{#1/#2}

% 2. Vertically stacked label (using a tight tabular for clean TikZ node rendering)
\newcommand{\stackedlabel}[2]{%
  \raisebox{-0.1ex}{\scriptsize\renewcommand{\arraystretch}{0.75}\begin{tabular}{@{}c@{}}#1\\#2\end{tabular}}%
}

\usepackage{expl3}

\ExplSyntaxOn
% Helper function to convert macro inputs (e.g. "7\s5 / 7\b13") into clean TikZ coordinate tags ("7s57b13")
\cs_new_protected:Npn \sanitize_tag:Nn #1#2
 {
   \str_set:Nx #1 { \detokenize{#2} }
   \str_remove_all:Nn #1 { ~ }
   \str_remove_all:Nn #1 { \c_backslash_str }
   \str_remove_all:Nn #1 { $ }
   \str_remove_all:Nn #1 { / }
   \str_remove_all:Nn #1 { ( }
   \str_remove_all:Nn #1 { ) }
   \str_remove_all:Nn #1 { + }
   \str_remove_all:Nn #1 { , }
   \str_remove_all:Nn #1 { \c_left_brace_str }
   \str_remove_all:Nn #1 { \c_right_brace_str }
 }

% Variable to hold current default muted strings
\newcommand{\defaultmute}{}
\newcommand{\setmute}[1]{\gdef\defaultmute{#1}}

% Optional: Wrapper to combine \subsubsection and \setmute in one command
\newcommand{\stringset}[2][]{%
  \setmute{#1}%
  \subsubsection{#2}%
}
\let\chordsubsection\stringset

% Helper to expand \defaultmute into actual numbers before calling \chordscheme
\newcommand{\renderchord}[1]{%
  \ifx\defaultmute\empty
    \chordscheme[#1]%
  \else
    \expanded{\noexpand\chordscheme[mute={\defaultmute}, \unexpanded{#1}]}%
  \fi
}

% Renewed \chordrow using \renderchord
\NewDocumentCommand{\chordrow}{m m m m m}{%
  \phantomsection%
  \noindent%
 \raisebox{-0.5\height}{\parbox[b]{1.8cm}{\raggedleft\large\bfseries #1}\hspace{0.1em}}%
  \raisebox{-0.5\height}{\renderchord{#2}}\hfill%
  \raisebox{-0.5\height}{\renderchord{#3}}\hfill%
  \raisebox{-0.5\height}{\renderchord{#4}}\hfill%
  \raisebox{-0.5\height}{\renderchord{#5}}%
  \sanitize_tag:Nn \l_tmpa_str {#1}%
  \tikz[remember~picture, overlay] \coordinate (s\arabic{section}-ss\arabic{subsection}-sss\arabic{subsubsection}-\l_tmpa_str);%
  \par\vspace{6pt}%
}

% Updated \drawbrace (automatically cleans its arguments)
\NewDocumentCommand{\drawbrace}{m m}{%
  \sanitize_tag:Nn \l_tmpa_str {#1}%
  \sanitize_tag:Nn \l_tmpb_str {#2}%
  \begin{tikzpicture}[remember~picture, overlay]
    \draw[decorate, decoration={brace, amplitude=5pt}, thick]
      ([xshift=10pt, yshift=0.5cm] s\arabic{section}-ss\arabic{subsection}-sss\arabic{subsubsection}-\l_tmpa_str) -- 
      ([xshift=10pt, yshift=-0.5cm] s\arabic{section}-ss\arabic{subsection}-sss\arabic{subsubsection}-\l_tmpa_str |- s\arabic{section}-ss\arabic{subsection}-sss\arabic{subsubsection}-\l_tmpb_str)
      node[midway, xshift=6pt, anchor=west] {};
  \end{tikzpicture}%
}
\ExplSyntaxOff
% ---------------------------------------------------------
% VANILLA GUITARCHORDSCHEMES SETTINGS
% ---------------------------------------------------------
\setchordscheme{
  chord-frets = 5,
  x-unit = 4.2mm,
  y-unit = 5.5mm,
  finger-format = \footnotesize\sffamily\bfseries,
  name-distance = -2mm,
  string-name-format = \nullfont,
}

% ---------------------------------------------------------
% HIERARCHY & TOC CONFIGURATION
% ---------------------------------------------------------
\setcounter{tocdepth}{3}
\setcounter{secnumdepth}{0}

\titleformat{\section}{\LARGE\bfseries}{\thesection}{0pt}{}
\titlespacing*{\section}{0pt}{12pt}{3pt}

\titleformat{\subsection}{\Large\bfseries}{\thesubsection}{0pt}{}
\titlespacing*{\subsection}{0pt}{8pt}{3pt}

\titleformat{\subsubsection}{\large\bfseries}{\thesubsubsection}{0pt}{}
\titlespacing*{\subsubsection}{0pt}{8pt}{0pt}
